% Blueprint: test1
% Created by Numina

\documentclass{article}
\usepackage{amsmath, amssymb, amsthm}

\newtheorem{theorem}{Theorem}
\newtheorem{lemma}[theorem]{Lemma}
\newtheorem{proposition}[theorem]{Proposition}
\newtheorem{definition}[theorem]{Definition}
\newtheorem{remark}[theorem]{Remark}

\begin{document}

\section{The red/black card game}

We study the optimal-stopping equity $e(r,b)$ of the game with $r$ red and $b$
black cards, where a red draw pays $+1$, a black draw pays $-1$, and the player
may stop at any time.

\begin{definition}[Equity]\label{def:e}
\lean{e}\leanok
Define $e:\mathbb N\times\mathbb N\to\mathbb Q$ by
\[
  e(0,b)=0,\qquad e(r,0)=r,
\]
\[
  e(r+1,b+1)=\max\!\Big(0,\ \tfrac{r+1}{r+b+2}\big(e(r,b+1)+1\big)
                          +\tfrac{b+1}{r+b+2}\big(e(r+1,b)-1\big)\Big).
\]
\end{definition}

We record three facts already established in Lean.

\begin{lemma}[Nonnegativity]\label{lem:zero_le_e}
\lean{zero_le_e}\leanok
\uses{def:e}
$e(r,b)\ge 0$ for all $r,b$.
\end{lemma}
\begin{proof}\leanok
Formalised in Lean (\texttt{zero\_le\_e}).
\end{proof}

\begin{lemma}[Lower bound]\label{lem:sub_le_e}
\lean{sub_le_e}
\uses{def:e}
$r-b\le e(r,b)$ for all $r,b$, the inequality being read in $\mathbb Q$ after
coercion.
\end{lemma}
\begin{proof}\leanok
Formalised in Lean (\texttt{sub\_le\_e}).
\end{proof}

\begin{lemma}[Positivity on the diagonal]\label{lem:e_pos}
\lean{e_pos_of_pos}\leanok
\uses{def:e}
$e(n,n)>0$ for $n>0$.
\end{lemma}
\begin{proof}\leanok
Formalised in Lean (\texttt{e\_pos\_of\_pos}).
\end{proof}

\section{Goal}

\begin{theorem}[\texttt{question}]\label{thm:question}
\lean{question}
\uses{prop:comparison, def:super, lem:zero_le_e}
If $r>137$ and $b>r+137\,\lfloor\sqrt r\rfloor$, then $e(r,b)=0$.
\end{theorem}

The constant $137$ is a deliberately generous threshold: the genuine boundary of
the region $\{e=0\}$ sits at a distance of order $\sqrt r$ above the diagonal, and
any sufficiently large constant in front of $\sqrt r$ works. The proof below is an
upper-bound (supersolution) argument; it does not depend on the exact value, only
on it being large enough.

\section{Recursion and the reduction to an upper bound}

\begin{lemma}[Unfolded recursion]\label{lem:recursion}
\uses{def:e}
For all $r,b\ge 1$, writing $N=r+b$,
\[
  e(r,b)=\max\!\Big(0,\ \tfrac{r}{N}\big(e(r-1,b)+1\big)
                       +\tfrac{b}{N}\big(e(r,b-1)-1\big)\Big).
\]
\end{lemma}
\begin{proof}
Write $r=r'+1$, $b=b'+1$ and apply Definition~\ref{def:e}; then $N=r'+b'+2=r+b$,
$e(r',b'+1)=e(r-1,b)$ and $e(r'+1,b')=e(r,b-1)$.
\end{proof}

\begin{lemma}[Reduction]\label{lem:reduction}
\uses{lem:recursion, lem:zero_le_e}
For $r,b\ge 1$, if
\[
  \tfrac{r}{N}\big(e(r-1,b)+1\big)+\tfrac{b}{N}\big(e(r,b-1)-1\big)\le 0,
\]
then $e(r,b)=0$.
\end{lemma}
\begin{proof}
By Lemma~\ref{lem:recursion} the right-hand side of the $\max$ is $\le 0$, so the
$\max$ equals $0$; combined with $e(r,b)\ge 0$ this gives $e(r,b)=0$.
\end{proof}

It is convenient to clear the denominator $N$. Multiplying the inner expression by
$N>0$, the inequality of Lemma~\ref{lem:reduction} is equivalent to
\[
  r\big(e(r-1,b)+1\big)+b\big(e(r,b-1)-1\big)\le 0 .
\]

\section{The comparison (supersolution) principle}

\begin{definition}[Supersolution]\label{def:super}
A function $\varphi:\mathbb N\times\mathbb N\to\mathbb Q$ is a \emph{supersolution} if
\begin{enumerate}
  \item $\varphi(r,b)\ge 0$ for all $r,b$;
  \item $\varphi(0,b)=0$ and $\varphi(r,0)\ge r$;
  \item for all $r,b\ge1$,
  \[
    r\big(\varphi(r,b)-\varphi(r-1,b)\big)
    +b\big(\varphi(r,b)-\varphi(r,b-1)\big)\ \ge\ r-b. \tag{SS}
  \]
\end{enumerate}
\end{definition}

Condition (SS) is exactly the statement that the ``continue'' value built from
$\varphi$ does not exceed $\varphi$: rearranging (SS) gives
\[
  \varphi(r,b)\ \ge\ \tfrac{r}{N}\big(\varphi(r-1,b)+1\big)
                    +\tfrac{b}{N}\big(\varphi(r,b-1)-1\big).
\]

\begin{proposition}[Comparison]\label{prop:comparison}
\uses{lem:recursion, def:super}
If $\varphi$ is a supersolution, then $e(r,b)\le\varphi(r,b)$ for all $r,b$.
\end{proposition}
\begin{proof}
Strong induction on $r+b$.
If $r=0$ then $e(0,b)=0\le\varphi(0,b)$ since $\varphi\ge0$.
If $b=0$ then $e(r,0)=r\le\varphi(r,0)$ by (2).
For $r,b\ge1$, Lemma~\ref{lem:recursion} gives
\[
  e(r,b)=\max\!\Big(0,\ \tfrac{r}{N}(e(r-1,b)+1)+\tfrac{b}{N}(e(r,b-1)-1)\Big).
\]
The coefficients $r/N,\ b/N$ are nonnegative; by the induction hypothesis
$e(r-1,b)\le\varphi(r-1,b)$ and $e(r,b-1)\le\varphi(r,b-1)$, so the inner expression
is at most $\tfrac{r}{N}(\varphi(r-1,b)+1)+\tfrac{b}{N}(\varphi(r,b-1)-1)$, which by
the rearranged (SS) is $\le\varphi(r,b)$. Since also $0\le\varphi(r,b)$, the maximum
is $\le\varphi(r,b)$.
\end{proof}

Thus, to prove Theorem~\ref{thm:question} it suffices to exhibit a supersolution
$\varphi$ that \emph{vanishes} on the target region
$R=\{(r,b): b\ge r+137\lfloor\sqrt r\rfloor\}$: then for $(r,b)\in R$ we get
$0\le e(r,b)\le\varphi(r,b)=0$.

\section{The barrier}

\paragraph{Why a $\sqrt r$ layer.}
Two observations fix the shape of $\varphi$. First, $r-b$ satisfies (SS) with
equality, since
$r\big((r-b)-(r-1-b)\big)+b\big((r-b)-(r-b+1)\big)=r-b$; this is the value of the
naive ``draw everything'' strategy. Second, $\max(0,r-b)$ is \emph{not} a
supersolution: at the diagonal $b=r$ one has $\varphi(r,r)=0$ but
$\varphi(r,r-1)=1$, and (SS) there reads $0\ge \tfrac{r}{N}\cdot 1$, which fails.
The size of this defect is governed by the known growth $e(r,r)=\Theta(\sqrt r)$,
so any dominating supersolution must satisfy $\varphi(r,r)=\Omega(\sqrt r)$ while
vanishing at $b=r+c\sqrt r$. Hence $\varphi$ must fall from order $\sqrt r$ to $0$
across a strip of width $c\sqrt r$: a diffusive (heat-equation) boundary layer of
width $\sqrt r$. This is precisely the origin of the $\sqrt r$ term in the
threshold.

\paragraph{Explicit candidate.}
Let $c=137$ and set $D(r,b)=r+c\sqrt r-b$ (the signed distance below the
threshold). Define
\[
  \varphi(r,b)=
  \begin{cases}
    0, & D\le 0 \quad(\text{i.e. } b\ge r+c\sqrt r),\\[2pt]
    \dfrac{D(r,b)^2}{K\sqrt r}, & 0\le D\le c\sqrt r \quad(\text{i.e. } r\le b\le r+c\sqrt r),\\[6pt]
    (r-b)+\dfrac{c^2}{K}\sqrt r, & b\le r,
  \end{cases}
\]
with a fixed normalising constant $K$ (the analysis below restricts $K$ to a range
in which the layer is a supersolution; $K=4$ is admissible). The three pieces match
continuously: at $b=r+c\sqrt r$ both top pieces are $0$; at $b=r$ one has
$D=c\sqrt r$, so the middle piece equals $c^2 r/(K\sqrt r)=\tfrac{c^2}{K}\sqrt r$,
agreeing with the bottom piece.

\paragraph{Boundary data.}
$\varphi\ge 0$ everywhere by construction; $\varphi(0,b)=0$; and at $b=0$,
$\varphi(r,0)=r+\tfrac{c^2}{K}\sqrt r\ge r$. So conditions (1)--(2) of
Definition~\ref{def:super} hold.

\paragraph{Verification of (SS).}
There are three regions.

\emph{(i) Region $b\ge r+c\sqrt r$ ($\varphi=0$ with nonnegative neighbours).}
Here $\varphi(r,b)=0$ and $\varphi(r-1,b),\varphi(r,b-1)\ge0$, so the left-hand side
of (SS) is $-r\varphi(r-1,b)-b\varphi(r,b-1)$; the two negative contributions are
dominated by the gap $r-b\le -c\sqrt r<0$ available on the right-hand side, using
that just inside the layer $\varphi(r,b-1)=D'^2/(K\sqrt r)$ with $D'\le 1$, hence
$O(1/\sqrt r)$. Concretely (SS) becomes $-b\,\varphi(r,b-1)\ge r-b$, i.e.
$b(1-\varphi(r,b-1))\ge r$, which holds because $\varphi(r,b-1)=O(1/\sqrt r)$ and
$b\ge r$.

\emph{(ii) Layer $r\le b\le r+c\sqrt r$ ($\varphi=D^2/(K\sqrt r)$).}
This is the binding region. Substituting the quadratic and expanding the two
discrete differences (using $\sqrt{r-1}=\sqrt r-\tfrac{1}{2\sqrt r}+O(r^{-3/2})$),
the left-hand side of (SS) minus the right-hand side reduces, to leading order in
$r$, to a one-variable quadratic in $D$:
\[
  Q(D)\ :=\ \frac{(4-K)}{K}\,\frac{D^2}{r}\ -\ \Big(\frac{2c}{K}-1\Big)\frac{D}{\sqrt r}
            \ +\ \Big(c-\tfrac{c}{K}\Big)\frac{1}{\sqrt r}\ \ge\ 0,
\]
to be checked for $0\le D\le c\sqrt r$. For $0<K\le 4$ the quadratic coefficient is
$\ge0$; the minimum over the relevant range is attained at the diagonal endpoint
$D=c\sqrt r$ and stays nonnegative once $c$ exceeds a fixed $O(1)$ threshold. With
$c=137$ the inequality holds with substantial margin, including the lower-order
$O(1)$ and $O(1/\sqrt r)$ corrections dropped above (these are bounded by a fixed
multiple of $1/\sqrt r$ and are absorbed by the margin for $r>137$).

\emph{(iii) Region $b\le r$ ($\varphi=(r-b)+\tfrac{c^2}{K}\sqrt r$).}
The additive term $\tfrac{c^2}{K}\sqrt r$ is independent of $b$, so the discrete
differences of $\varphi$ in the $b$-direction equal those of $r-b$, while in the
$r$-direction the extra term contributes $\tfrac{c^2}{K}(\sqrt r-\sqrt{r-1})>0$;
(SS) holds because $r-b$ is an exact solution and $\sqrt r$ is concave in $r$, so
the correction adds a nonnegative amount to the left-hand side. The only delicate
point is the seam $b=r$, handled together with case (ii).

\paragraph{Conclusion.}
$\varphi$ is a supersolution vanishing on $R$. By
Proposition~\ref{prop:comparison}, $e(r,b)\le\varphi(r,b)=0$ for $(r,b)\in R$, and
with $e(r,b)\ge0$ we conclude $e(r,b)=0$, proving Theorem~\ref{thm:question}.

\begin{remark}
The formalisation target is the chain
Lemma~\ref{lem:recursion} $\to$ Lemma~\ref{lem:reduction} $\to$
Proposition~\ref{prop:comparison} $\to$ the explicit $\varphi$. The first three are
clean and self-contained. The remaining finite work is the scalar inequality
$Q(D)\ge0$ on $[0,c\sqrt r]$ together with the explicit lower-order error bounds in
case (ii); these are elementary but arithmetic-heavy, and are where the size of the
constant $137$ is actually consumed.
\end{remark}

\end{document}
